\documentclass[12pt,a4paper]{article}

% ── Idioma y codificación ─────────────────────────────────────────────────────
% "english" se carga como idioma secundario para el abstract en inglés
% (obligatorio segun la guia de Scientia Agropecuaria); "spanish" al final
% queda como idioma principal/por defecto del documento.
\usepackage[english,spanish,es-noshorthands]{babel}
\AtBeginDocument{\spanishplainpercent}
\usepackage[T1]{fontenc}
\addto\captionsspanish{\renewcommand*{\tablename}{Tabla}}
\AtBeginDocument{\renewcommand{\appendixname}{Appendix}}

% Etiqueta segun el idioma activo de babel: "Keywords" dentro del bloque
% \begin{otherlanguage}{english}...\end{otherlanguage} del abstract en
% ingles, "Palabras clave" en el resto del documento (espanol).
\providecommand{\keywords}[1]{\noindent\textbf{%
    \iflanguage{english}{Keywords}{Palabras clave}: } #1}

% ── ORCID ─────────────────────────────────────────────────────────────────────
\usepackage{orcidlink}

% ── Interlineado a doble espacio (Scientia Agropecuaria exige doble espacio) ──
\usepackage[doublespacing]{setspace}
\usepackage{microtype}

% ── Numeración de líneas: DESACTIVADA ────────────────────────────────────────
% Se conservan \linenumbers / \nolinenumbers como órdenes vacías para no tocar
% los archivos 03/04/05/06. Para reactivarla: \usepackage{lineno} y borrar las
% dos líneas siguientes.
\newcommand{\linenumbers}{}
\newcommand{\nolinenumbers}{}

% ── Diseño de página ──────────────────────────────────────────────────────────
\usepackage{multicol}
\usepackage{float}
\usepackage{enumitem}
\usepackage{emptypage}
\usepackage{csquotes}
\usepackage{geometry}

\geometry{
    a4paper,
    margin    = 25mm,
    top       = 25mm,
    right     = 25mm,
    left      = 25mm,
    bottom    = 25mm,
    columnsep = 1cm,
}

% ── Manuscrito a 1 columna (Scientia Agropecuaria: formato de envío) ─────────
% El cuerpo (04_body.tex) usa \begin{multicols}{2}...\end{multicols} para
% intercalar tablas/figuras en modo landscape entre bloques de texto a dos
% columnas. En vez de reescribir cada aparición, se redefine el entorno
% "multicols" para que se comporte como una columna simple (equivalente a
% \begin{multicols}{1}): el maquetado final a 2 columnas, si la revista lo
% aplica, lo hace la editorial tras la aceptación.
\let\multicols\relax
\let\endmulticols\relax
\newenvironment{multicols}[1]{}{}

% ── Tolerancias de composición ────────────────────────────────────────────────
\setlength{\emergencystretch}{3em}
\hbadness=10000
\vbadness=10000
\tolerance=2000
\overfullrule=0pt
\hfuzz=18pt

% ── Interacciones lineno con otros entornos ───────────────────────────────────
\AtBeginEnvironment{longtblr}{\nolinenumbers\hfuzz=18pt}
\AtBeginEnvironment{tblr}{\nolinenumbers}
\AtEndEnvironment{tblr}{\linenumbers}
\AtBeginEnvironment{multicols}{\nolinenumbers}
\AtEndEnvironment{multicols}{\linenumbers}
\AddToHook{env/landscape/begin}[lineno-fix]{\nolinenumbers}
\AddToHook{env/landscape/end}[lineno-fix]{\linenumbers}
\DeclareHookRule{env/landscape/begin}{lineno-fix}{before}{pdflscape}

% ── Gráficos y figuras ────────────────────────────────────────────────────────
\usepackage{graphicx}
\usepackage{caption}
\usepackage{subcaption}
\usepackage{tabularray}
\graphicspath{{figures/}{./figures/}}
\renewcommand{\thefigure}{\arabic{figure}}

% ── Diagramas esquemáticos propios (TikZ) ─────────────────────────────────────
\usepackage{tikz}
\usetikzlibrary{arrows.meta,positioning,fit,backgrounds,shapes.geometric}

% ── Tablas ────────────────────────────────────────────────────────────────────
\usepackage{tabularx}
\usepackage{longtable}
\usepackage{multirow}
\usepackage{booktabs}
\usepackage{array}
\usepackage{rotating}

\UseTblrLibrary{booktabs}
\UseTblrLibrary{siunitx}

% ── Estilo de leyendas ────────────────────────────────────────────────────────
\captionsetup[table]{
    skip=6pt,
    font={small,sf,stretch=1.5},
    labelfont={bf},
    textfont={it},
    labelsep=newline,
    justification=raggedright,
    singlelinecheck=false,
}
\captionsetup[figure]{
    skip=6pt,
    font={small,sf,stretch=1.5},
    labelfont={bf},
    textfont={it},
    labelsep=newline,
    justification=justified,
    singlelinecheck=false,
}

% ── Entorno Nota (pie de tabla/figura) ────────────────────────────────────────
\NewDocumentEnvironment{Nota}{O{Nota} m b}{%
    \nolinenumbers
    \par\vspace{10pt plus2pt minus2pt}%
    \begin{minipage}{#2}%
        \nolinenumbers
        \onehalfspacing\sffamily\small\textit{\textbf{#1}.}\ #3%
    \end{minipage}%
}{}

% ── Secciones (formato Elsevier) ──────────────────────────────────────────────
\usepackage{titlesec}
\titleformat{\section}{\normalsize\bfseries}{\thesection}{1em}{}
\titleformat{\subsection}{\normalsize\bfseries}{\thesubsection}{1em}{}
\titleformat{\subsubsection}{\normalsize\itshape}{\thesubsubsection}{1em}{}
\titleformat{\paragraph}[runin]{\normalsize\bfseries}{}{0pt}{}[.]

% ── Miscelánea ────────────────────────────────────────────────────────────────
\usepackage{pdflscape}
\usepackage{setspace}
\usepackage{pdfpages}
\usepackage[Export]{adjustbox}
\usepackage{url}
\usepackage{doi}
\usepackage{chngcntr}
\usepackage{soul}
\usepackage{xcolor}

% ── Matemáticas ───────────────────────────────────────────────────────────────
\usepackage{amsmath}
\usepackage{amsfonts}
\usepackage{mathtools}
\usepackage{mathrsfs}
\usepackage{amssymb}

% ── Tipografía: Times New Roman (Scientia Agropecuaria exige Times 12pt) ──────
% Debe cargarse DESPUES de amsmath/amsfonts/mathtools/amssymb (orden exigido
% por la documentacion de newtx) para evitar conflictos de simbolos
% duplicados (p. ej. "\Bbbk already defined").
\usepackage{newtxtext}
\usepackage{newtxmath}

\usepackage{siunitx}
\sisetup{
    separate-uncertainty = true,
    multi-part-units     = repeat,
}
\NewDocumentCommand\Laplace{}{\mathscr{L}}
\renewcommand{\theequation}{\arabic{equation}}

% ── Código fuente ─────────────────────────────────────────────────────────────
\usepackage{listings}
\definecolor{codegray}{rgb}{0.95,0.95,0.95}
\definecolor{codeblue}{rgb}{0.15,0.30,0.65}
\definecolor{codered}{rgb}{0.65,0.10,0.10}
\definecolor{codegreen}{rgb}{0.10,0.50,0.20}
\lstset{
    language={},
    keywordstyle=\color{codeblue}\bfseries,
    stringstyle=\color{codered},
    commentstyle=\color{codegreen}\itshape,
    basicstyle=\ttfamily\footnotesize,
    showstringspaces=false,
    numbers=left,
    numberstyle=\tiny\color{gray},
    breaklines=true,
    breakatwhitespace=false,
    postbreak=\mbox{\textcolor{gray}{$\hookrightarrow$}\space},
    frame=single,
    backgroundcolor=\color{codegray},
    captionpos=b,
    xleftmargin=1.5em,
    columns=flexible,
    tabsize=2,
}

% ── Comando para título del manuscrito ────────────────────────────────────────
\newcommand{\tituloprincipal}[2]{%
    \begin{center}
        {\fontsize{14}{17}\selectfont\textbf{#1}\\[0.5em]}
        {\fontsize{11}{14}\selectfont #2}
    \end{center}
}

% ── Comando para métricas inline ──────────────────────────────────────────────
\newcommand{\metric}[3]{$#1 = \SI{#2}{#3}$}

% ── Hyperref ──────────────────────────────────────────────────────────────────
\usepackage{hyperref}
\hypersetup{
	colorlinks   = true,
	citecolor    = blue!75!black,
	linkcolor    = blue!60!black,
	urlcolor     = blue!75!black,
	filecolor    = blue!75!black,
	pdftitle     = {Aprendizaje automatico y control avanzado para la
		regulacion del pH en la industria azucarera: una revision de
		alcance y analisis bibliometrico},
	pdfauthor    = {Abraham Castillo},
	pdfsubject   = {Revision de alcance PRISMA-ScR y analisis bibliometrico
		sobre aprendizaje automatico y control avanzado del pH en la
		industria azucarera},
	pdfkeywords  = {control de pH; industria azucarera; aprendizaje automatico;
		control predictivo; revision de alcance; analisis bibliometrico;
		brecha de investigacion},
	pdfstartview = {FitH},
}

% ── Biblatex / APA (Scientia Agropecuaria exige APA 6.a ed. o posterior) ──────
\usepackage[
backend    = biber,
style      = apa,
sorting    = nyt,
hyperref   = auto,
]{biblatex}
\DeclareLanguageMapping{spanish}{spanish-apa}
\addbibresource{bibliography/pHazucar.bib}


\let\origparencite\parencite
\RenewDocumentCommand{\parencite}{O{} O{} m}{\textcolor{blue!75!black}{\origparencite[#1][#2]{#3}}}
\let\origtextcite\textcite
\RenewDocumentCommand{\textcite}{O{} O{} m}{\textcolor{blue!75!black}{\origtextcite[#1][#2]{#3}}}


% ── Profundidad de numeración ─────────────────────────────────────────────────
\setcounter{secnumdepth}{4}
\setcounter{tocdepth}{3}

% ── Datos pendientes (se rellenan en un único archivo) ───────────────────────
% 00_datos_pendientes.tex — ÚNICO lugar donde se completan los datos que aún no
% están en el repositorio. Mientras una macro conserve su definición por
% defecto, se imprime en ROJO entre corchetes (visible en el PDF).
% Para completarla basta redefinirla, p. ej.:
%     \renewcommand{\KappaInter}{0.85}
% y el valor se propaga al resumen, a los resultados, a las tablas y a los
% apéndices. NO escribir cifras directamente en 02/04/06.
%
% Dependencias entre cifras (comprobar al rellenar):
%   * SciELO y LILACS: 0 resultados (confirmado por el autor); el embudo
%     541→353→239→230 y la Figura PRISMA no cambian.
%   * El Revisor 2 cribó los 230 registros (no una muestra): \NPositivosMuestra=2.
%     Fuente: 03_filtered_data/4.filtrado_ollama.xlsx (decision_revisor1/2),
%     script 05_kappa_cohen.py. Wilson 95 %: 2/2 -> [34.2, 100]; 230/230 -> [98.4, 100].

\newcommand{\DatoPendiente}[1]{\textbf{\textcolor{red}{[#1]}}}

% ── Búsqueda en SciELO y LILACS (Tabla A1) ───────────────────────────────────
\newcommand{\NSciELO}{0}
\newcommand{\NLILACS}{0}
\newcommand{\NSciELOLILACS}{0}
\newcommand{\NTotalAmpliado}{541}
\newcommand{\NUnicosAmpliado}{353}

% ── Concordancia inter-humano (métrica primaria; 230 registros) ──────────────
\newcommand{\KappaInter}{1.000}
\newcommand{\PzeroInter}{100\,\%}
\newcommand{\PABAKInter}{1.000}
\newcommand{\SensInter}{100\,\%}
\newcommand{\PrecInter}{100\,\%}
\newcommand{\FUnoInter}{1.000}
\newcommand{\ICInter}{[34.2; 100]\,\% para la sensibilidad y la precisión (2/2) y [98.4; 100]\,\% para $P_0$ (230/230)}
\newcommand{\PrevInter}{0.87\,\%}
\newcommand{\NPositivosMuestra}{2}
% Matriz de confusión inter-humano
\newcommand{\VPInter}{2}
\newcommand{\FNInter}{0}
\newcommand{\FPInter}{0}
\newcommand{\VNInter}{228}
\newcommand{\TotUnoI}{2}
\newcommand{\TotUnoE}{228}
\newcommand{\TotDosI}{2}
\newcommand{\TotDosE}{228}

% ── Enlaces de datos y protocolo ─────────────────────────────────────────────
\newcommand{\URLOSF}{https://doi.org/10.17605/OSF.IO/4CH35}
% Repositorio confirmado por el autor (HTTP 200 el 24-oct-2026).
\newcommand{\URLGitHub}{https://github.com/acastillo1975/rsl-ph-azucarera}
\newcommand{\OSFTexto}{doi:10.17605/OSF.IO/4CH35}
\newcommand{\GitHubTexto}{github.com/acastillo1975/rsl-ph-azucarera}

% ── Registro retrospectivo en OSF (completar al registrar) ──────────────────
% Fecha en que se envió el registro en OSF (posterior al 22-jun-2026).
\newcommand{\FechaRegistroOSF}{24~de octubre de 2026}

% ── Fecha de ejecución de la búsqueda ────────────────────────────────────────
% Confirmada por el autor: 22-jun-2026 (los RIS de Scopus llevan «Export Date:
% 23 June 2026», un día después de ejecutar la búsqueda).
\newcommand{\FechaBusqueda}{22~de junio de 2026}

% ── Co-revisor ───────────────────────────────────────────────────────────────
\newcommand{\NombreCoRevisor}{D.~Barreto}

\endinput


\endinput
